\section{Further fixed-coset transfer arguments}
\label{l:transfer-supplement}

The main transfer assigns each joint restriction type to one small
factor, or covers the representation by small isotypic projections.
The alternatives here retain different information. The first two
arguments average many low-type projections before recovering a test
vector; the third counts whole joint types. We also record parameter
choices for these and the main estimates. Finally, a separate character
construction removes whole longest lines instead of assigning individual
boxes to rows or columns. Its dimension ratio, reciprocal-degree proof,
and transfer are kept together.

For the projection arguments and substitutions below, use the
fixed-coset data of Section~\ref{l:comparisons-section}. Namely, let
$G=S_n$, partition $[n]=M_1\sqcup\cdots\sqcup M_b$ into nonempty
blocks, and let $H_i=\Sym(M_i)$ fix the complement. Let $f\ge0$ have
mass at most one and satisfy
\[
 f(gH_i)\le \frac{B}{[G:H_i]}\qquad(g\in G,\ 1\le i\le b).
\]
Write $D=\dim\rho$ for the degree of the ambient irreducible and
$C_u=\sup_{m\ge1}\sum_{\lambda\vdash m}D_\lambda^{-u}$.
The peeling construction concerns a single symmetric group first;
its final transfer states its own equal-block data.

\subsection{Averages of many low-type projections}

\begin{proposition}[Coefficient-space average]\label{l:coefficient-average}
For $b=32$ and cap $B=4$,
\[
 \|\widehat f(\rho)\|_{\op}\le\frac83\sqrt{C_1}\,D^{-5/16}.
\]
\end{proposition}
\begin{proof}
Let $P_i$ now select factor degrees at most $D^{1/8}$. At most eight factors of any product type can exceed that threshold, so $\bar P=32^{-1}\sum_iP_i\succeq3I/4$. Lemma~\ref{l:coefficient-evaluation}, the coset cap, and full-group Schur orthogonality give
\[
 |\langle w,\widehat f(\rho)P_iv\rangle|^2
 \le4C_1D^{-5/8}\|w\|^2\|P_iv\|^2.
\]
Here the selected coefficient space has dimension at most $C_1D^{3/8}$. Thus $\|\widehat fP_i\|_{\op}\le2\sqrt{C_1}D^{-5/16}$. Average over $i$ and multiply by $\bar P^{-1}$, whose norm is at most $4/3$.
\end{proof}

\begin{proposition}[Orbit-space average]\label{l:orbit-average}
For $b=256$ and cap $B=4$,
\[
 \|\widehat f(\rho)\|_{\op}\le4\sqrt{C_{32}}D^{-1/4}.
\]
\end{proposition}
\begin{proof}
Let $P_i$ select degrees at most $D^{2/b}$. At least half the factors of every product type meet this condition, so $\bar P=b^{-1}\sum_iP_i\succeq I/2$. A vector in the range of $P_i$ has orbit span of dimension at most $C_{32}D^{68/b}$ by Lemma~\ref{l:single-orbit}. The coset cap and Schur averaging, exactly as in Proposition~\ref{l:orbit-span}, give
\[
 \|\widehat fP_i\|_{\op}\le2\sqrt{C_{32}}D^{-1/2+34/b}
 \le2\sqrt{C_{32}}D^{-1/4},
\]
since $68/256\le1/2$. Average and use $\|\bar P^{-1}\|_{\op}\le2$. This proof retains the geometric averaging mechanism and only uses the reciprocal power $32$.
\end{proof}

\begin{proposition}[Four-factor coefficient decomposition]\label{l:four-coefficient}
If there are four blocks and the subprobability cap is $B=2$, then
\[
 \|\widehat f(\rho)\|_{\op}\le2C_{1/4}^2D^{-1/8}.
\]
\end{proposition}
\begin{proof}
Decompose a vector orthogonally into full product types $v=\sum_{\boldsymbol\tau}v_{\boldsymbol\tau}$. If $d_i$ are the four factor degrees, their product is at most $D$. The sum of $(\prod_i d_i)^{-1/4}$ over all possible product types is at most $C_{1/4}^4$. Thus the number $T$ of occurring types is at most $C_{1/4}^4D^{1/4}$.

For one type choose $i$ with $d_i\le D^{1/4}$. The coefficient function on $gH_i$ belongs to a single type's coefficient space, even when that type occurs repeatedly. Lemma~\ref{l:coefficient-evaluation} bounds its supremum by $d_i$ times its uniform root mean square. Sum over cosets with cap two, then apply Cauchy--Schwarz over the uniform cosets. Full-group Schur orthogonality yields
\[
 \sum_gf(g)|\langle w,\rho(g)v_{\boldsymbol\tau}\rangle|
 \le2d_iD^{-1/2}\|w\|\|v_{\boldsymbol\tau}\|
 \le2D^{-1/4}\|w\|\|v_{\boldsymbol\tau}\|.
\]
Finally $\sum_{\boldsymbol\tau}\|v_{\boldsymbol\tau}\|\le\sqrt T\|v\|$, giving the claimed exponent and constant. The count here is of entire product types, unlike the blockwise assignment in Proposition~\ref{l:coefficient-operator}.
\end{proof}

\subsection{Exact substitutions}
\label{l:exact-substitutions}

For applications using a prescribed block count or reciprocal-degree power, the following substitutions give explicit constants and exponents. They are convenient choices, not optimal convolution counts. In Proposition~\ref{l:orbit-span}, $u=16$ gives exponent $-(1-18/b)/2$; $u=2$ gives $-(1-4/b)/2$; $u=1$, $b=8$ gives $-5/16$. Each substitution changes only the inequality $\sum_{a\le D^{1/b}}a^2\le C_uD^{(u+2)/b}$.

For Hilbert--Schmidt bounds, the coefficient proof with $u=2$, $b=16$, $B=3$ gives $D\|\widehat f\|_{\HS}^2\le9bC_2D^{4/b}$. The pointwise character proof with $u=2$, $b=8$, $B=1$ gives $\|\widehat f\|_{\HS}^2\le bC_2D^{-1+6/b}$. The central-isotypic proof with $u=16$, $b=64$, $B=4$ gives $4bC_{16}D^{-1+18/b}$. With respectively four, sixteen, and thirty-two convolutions, the nonsign exponents are $-2$, $-3$, and $-22$. The coefficient proof with $u=10$ gives $B^2bC_{10}D^{-1+12/b}$; its sharper branching alternative is Proposition~\ref{l:branching-hs}. These substitutions preserve the distinct proofs of the estimates even when a stronger estimate is also available.

\subsection{A separate construction by peeling longest lines}

The next construction chooses whole lines successively, instead of
assigning individual boxes.  Its index estimate is weaker, but its
stepwise dimension inequality gives another proof of inverse-degree
summability.  The occurrence argument uses the horizontal- and
vertical-strip cases of the Pieri rule \cite{sagan}.

\begin{lemma}[A subcharacter obtained by peeling]
\label{l:peeling-character}
Put $C_{\rm peel}=512^2$.  Every irreducible representation of $S_m$
of degree $D$ has a one-dimensional character in its restriction to
some subgroup $J$ with $[S_m:J]\le D^{C_{\rm peel}}$.
At a step removing a longest first row or first column of length $r$,
leaving size $s$ and degree $D'$, one has
\begin{equation}\label{l:peeling-ratio}
 \frac D{D'}\ge\binom{r+s}r^{1/C_{\rm peel}}.
\end{equation}
The degree of the empty remainder is one.
\end{lemma}
\begin{proof}
Remove the longer of the first row and first column of the current
diagram.  In a row removal, the realigned lower diagram $\mu$ is
contained in the original diagram $\lambda$, and $\lambda/\mu$
is a horizontal strip: every column loses exactly its bottom box.
The Pieri rule and Frobenius reciprocity show that restriction to
$S_s\times S_r$ contains $V_\mu$ tensored with the trivial
representation of $S_r$.  For a column removal use the sign
representation and the transposed argument.  Iteration supplies a
product subgroup and a product of trivial and sign characters
occurring in the original representation.

To prove the index bound, it remains to compare dimensions at each
step.  Transpose so that a row is removed; then the diagram has height
at most $r$.  If $c_j$ counts its lower boxes in column $j$, the
hook formula gives
\[
 \frac D{D'}=\frac{\binom{r+s}r}{R},\qquad
 R=\prod_{j=1}^r\left(1+\frac{c_j}{r-j+1}\right),
 \qquad \sum_jc_j=s.
\]
The case $s=0$ is immediate.  For $r\ge512$ and $s\ge16r$,
the bound $c_j\le r$ gives
\[
 R\le\binom{2r}r\le4^r,
 \qquad \binom{r+s}r\ge(1+s/r)^r\ge17^r.
\]
Thus $R\le\binom{r+s}r^{1/2}$.  For $1\le s\le16r$,
the opposite ordering of $c_j$ and $(r-j+1)^{-1}$ instead gives
\[
 \log R\le\frac sr\sum_{j=1}^r\frac1j
       \le\frac{s(1+\log r)}r,
 \qquad \log\binom{r+s}r\ge s\log(1+r/s).
\]
Since $(1+\log r)/r$ decreases for $r\ge512$ and is at most
$\frac12\log(17/16)$ there, this also gives
$D/D'\ge\binom{r+s}r^{1/2}$.

If $r<512$ and $s>0$, then $r\ge2$ and $r+s\le r^2$.
Each standard ordering of the lower diagram yields two distinct
standard tableaux: fill the first row before that ordering, or fill
its first $r-1$ boxes, the first lower box, the last top-row box,
and then the remaining lower boxes in that ordering.  The first lower
box is in column one, so the second filling is also valid.  These
two constructions are injective and have disjoint images.  Hence
$D/D'\ge2$, while
$\binom{r+s}r\le2^{r^2}\le2^{C_{\rm peel}}$.
This proves \eqref{l:peeling-ratio} in every case.  Multiplying it
through the successive steps telescopes the dimension ratios.  The
product of the binomial coefficients is exactly the index of the
product subgroup constructed above, proving the assertion.
\end{proof}

\begin{lemma}[Inverse-first-degree summability from peeling]
\label{l:peeling-inverse-one}
The preceding stepwise estimate independently implies
\[
 \sum_{\lambda\vdash m:\,D_\lambda>1}D_\lambda^{-1}\longrightarrow0.
\]
Consequently, for all sufficiently large $m$ and every $Y\ge1$,
$S_m$ has at most $3Y$ irreducible types of degree at most $Y$.
\end{lemma}
\begin{proof}
We use $p(h)\le e^{3\sqrt h}$, which follows by evaluating the
partition generating function at $e^{-1/\sqrt h}$: the logarithm
of the product is at most
$\sqrt h\sum_{j\ge1}j^{-2}\le2\sqrt h$, and the coefficient
bound adds $\sqrt h$.

The first removed length $r$ is at least $\sqrt m$ because both
width and height are at most $r$.  If $r<m/2$,
\eqref{l:peeling-ratio} gives
\[
 \log D_\lambda\ge
       \frac{r\log(m/r)}{C_{\rm peel}}.
\]
The concavity of $r\log(m/r)$ bounds it below on
$[\sqrt m,m/2]$ by the smaller of
$\frac12\sqrt m\log m$ and $\frac12m\log2$.
After division by $\sqrt m$, this lower bound tends to infinity.
It therefore dominates the logarithm $3\sqrt m$ of the number of
all diagrams, proving a vanishing inverse sum in this range.

If $m/2\le r<m$, put $h=m-r\ge1$.  At most $2p(h)$ diagrams
have this value, counting both orientations, and
\eqref{l:peeling-ratio} gives
$D_\lambda\ge(m/h)^{h/C_{\rm peel}}$.
For fixed $h$ their inverse-degree contribution tends to zero.  It
is bounded by
$2e^{3\sqrt h}2^{-h/C_{\rm peel}}$, a summable sequence.
Dominated convergence handles this range.  The remaining diagrams
are the row and column, of degree one.  Once the nonlinear reciprocal
sum is at most one, at most $Y$ nonlinear types have degree at most
$Y$; adding the two linear types gives $Y+2\le3Y$.
\end{proof}

\begin{corollary}[The peeling character transfer]
\label{l:peeling-lift}
Let $b=2^{22}$ and partition $n=bm$ points into $b$ sets of size
$m$.  Let $H_i$ permute the $i$th set.  If $f\ge0$ has mass at
most one and $f(gH_i)\le2/[S_n:H_i]$ for every $g,i$, then for
all sufficiently large $m$ and every irreducible of degree $D$,
\begin{equation}\label{l:peeling-lift-bound}
 \|\widehat f(\rho)\|_{\HS}^2
 \le6bD^{-1+(C_{\rm peel}+2)/b}.
\end{equation}
\end{corollary}
\begin{proof}
Lemma~\ref{l:subcharacter-isotypic}, using
Lemma~\ref{l:peeling-character}, bounds a full $H_i$-isotypic
projection of degree $d_\tau$ by
$2d_\tau^{C_{\rm peel}+1}/D$ in squared Hilbert--Schmidt norm.
As in Theorem~\ref{l:character-lift}, the projections with
$d_\tau\le Y=D^{1/b}$ cover the identity.  By
Lemma~\ref{l:peeling-inverse-one} there are at most $3Y$ such
types for each factor.  Summing gives
$D^{-1}\cdot2b\cdot3Y\cdot Y^{C_{\rm peel}+1}$,
which is \eqref{l:peeling-lift-bound}.
\end{proof}

Here $(C_{\rm peel}+2)/b=(262144+2)/4194304<1/8$.
Thus four convolutions again tend to uniform whenever the retained
mass tends to one and the original sign mean tends to zero, by
Proposition~\ref{l:amplification}.  The peeling and box-assignment
proofs supply the subcharacter hypothesis separately; neither uses
the other's subgroup construction.
